\begin{figure}
  \centering
  \begin{tikzpicture}[x=1cm,y=1cm]
    \draw[ham/brace] (0,.62)--(3,.62);
    \node[font=\small,anchor=south] at (1.5,.77) {$h_j$};
    \foreach \hamGridRow/\hamPanelShift/\hamGridName in {0/0/{0},1/1/{1/3},2/2/{2/3}} {
    \draw[hamblue] (0,-\hamGridRow)--(11.4,-\hamGridRow);
    \foreach \hamGridIndex in {0,1,2,3} {
    \pgfmathsetmacro{\hamGridBoundary}{3*\hamGridIndex+\hamPanelShift}
    \draw (\hamGridBoundary,-\hamGridRow-.12)--(\hamGridBoundary,-\hamGridRow+.12);
    }
    \node[anchor=east,font=\small] at (-.25,-\hamGridRow) {$t=\hamGridName$};
    }
    \draw[hamteal,line width=3pt] (1,-1)--(4,-1);
    \draw[hamteal] (1,-1.2)--(1,-.8);
    \draw[hamteal] (4,-1.2)--(4,-.8);
    \draw[dashed] (2.8,.2)--(2.8,-2.55);
    \draw[dashed] (3.4,.2)--(3.4,-2.55);
    \node[anchor=north,font=\small] at (2.8,-2.6) {$s$};
    \node[anchor=north,font=\small] at (3.4,-2.6) {$r$};
    \node[ham/note,text width=5.4cm] at (7.4,-2.95)
    {$3|s-r|<h_j\le6|s-r|$\\both points in one interval};
  \end{tikzpicture}
  \caption[Three shifted interval grids]{The three grids at one selected scale.
    Their boundaries are separated by $h_j/3>|s-r|$, so at least one grid has no
    boundary between $s$ and $r$. The full nested construction uses endpoints
    $h_j(n+(-1)^jt)$ for $t=0,1/3,2/3$; the shift alternates with the level. The
    resulting infinite sum dominates $|s-r|^{-a}$ off the diagonal. This is not
    a claim that an arbitrary finite partial sum already dominates it. See
    \ProofRef{Theorem}{T:full-upper}.}
  \label{Fig:dyadic-cover}
\end{figure}
